\documentclass[a4paper,10pt]{article}
\usepackage[utf8]{inputenc}
\usepackage[T1]{fontenc}
\usepackage[spanish]{babel}
\usepackage{geometry}
\usepackage{array}
\usepackage{colortbl}
\usepackage{xcolor}
\usepackage{booktabs}
\usepackage{enumitem}
\usepackage{fancyhdr}
\usepackage{tikz}
\usepackage{tcolorbox}
\usepackage{longtable}
\usepackage{multicol}

\tcbuselibrary{skins,breakable}

\geometry{left=2cm, right=2cm, top=2.5cm, bottom=2cm}

% Colores
\definecolor{purpura}{HTML}{8E44AD}
\definecolor{purpuraclaro}{HTML}{E8D5F5}
\definecolor{azulg}{HTML}{3498DB}
\definecolor{rojog}{HTML}{E74C3C}
\definecolor{gris}{HTML}{F2F2F2}
\definecolor{verde}{HTML}{27AE60}
\definecolor{naranja}{HTML}{E67E22}
\definecolor{dorado}{HTML}{D4AC0D}

\pagestyle{fancy}
\fancyhf{}
\fancyhead[L]{\small\textbf{Nuevo Compa\~{n}eros 1 --- U03: La Familia}}
\fancyhead[R]{\small\textbf{P\'ildora formativa 3.4 --- Posesivos}}
\fancyfoot[C]{\thepage}

\setlength{\parindent}{0pt}
\renewcommand{\arraystretch}{1.3}

% Estilo de tcolorbox para cabeceras de diapositiva
\newtcolorbox{diapoheader}[1]{
  colback=purpura,
  colframe=purpura,
  coltext=white,
  fontupper=\large\bfseries,
  boxrule=0pt,
  arc=3pt,
  left=8pt, right=8pt, top=5pt, bottom=5pt,
  title={},
  before upper={#1}
}

% Estilo para fase MARS
\newtcolorbox{marsfase}[1]{
  colback=purpuraclaro,
  colframe=purpura,
  coltext=purpura,
  fontupper=\small\bfseries,
  boxrule=0.8pt,
  arc=2pt,
  left=5pt, right=5pt, top=2pt, bottom=2pt,
  width=0.48\textwidth,
  title={},
  before upper={#1}
}

% Estilo para preguntas de awareness
\newtcolorbox{awarenessbox}{
  colback=purpuraclaro,
  colframe=purpura,
  boxrule=0.8pt,
  arc=3pt,
  left=8pt, right=8pt, top=5pt, bottom=5pt,
}

% Estilo para instrucciones de juego
\newtcolorbox{gamebox}{
  colback=dorado!10,
  colframe=dorado!80!black,
  boxrule=0.8pt,
  arc=3pt,
  left=8pt, right=8pt, top=5pt, bottom=5pt,
}

% Estilo para concepto clave
\newtcolorbox{conceptbox}{
  colback=purpura!8,
  colframe=purpura,
  boxrule=1.2pt,
  arc=4pt,
  left=10pt, right=10pt, top=8pt, bottom=8pt,
}

% Comando para posesivos en color
\newcommand{\pos}[1]{\textbf{\textcolor{purpura}{#1}}}

\begin{document}

%% ============================================================
%% PÁGINA 1: PORTADA + CONTENIDO PARA EL PROFESOR
%% ============================================================

\begin{center}
{\LARGE\textbf{\textcolor{purpura}{P\'ILDORA FORMATIVA 3.4}}}\\[0.3cm]
{\Large\textbf{POSESIVOS: CONCORDANCIA CON EL OBJETO POSE\'IDO,}}\\[0.1cm]
{\Large\textbf{NO CON EL POSEEDOR}}\\[0.3cm]
{\normalsize Secci\'on Gram\'atica (p.37) --- Proyectable --- 4 diapositivas}\\[0.1cm]
{\small Enfoque inductivo: el estudiante deduce, el profesor no enuncia la regla.}
\end{center}

\vspace{0.5cm}

%% --- Resumen de diapositivas ---

\begin{tcolorbox}[colback=gris, colframe=purpura, title={\textbf{Resumen de diapositivas}}, fonttitle=\bfseries\color{white}, colbacktitle=purpura, arc=3pt]
\renewcommand{\arraystretch}{1.5}
\begin{tabular}{c >{\raggedright\arraybackslash}p{5.5cm} >{\raggedright\arraybackslash}p{6.5cm}}
\textbf{Diap.} & \textbf{T\'itulo} & \textbf{Fase MARS} \\
\midrule
1 & Las flechas de posesi\'on & M --- Modelling \\
2 & ?`Qu\'e cambia a la vez? & A --- Awareness Raising \\
3 & ?`Todos funcionan igual? + Desaf\'io Rel\'ampago & A --- Awareness + S --- Structured Production \\
4 & Comprobad con el libro --- Actividad 5 & R --- Receptive Processing + S --- Structured Production \\
\end{tabular}
\end{tcolorbox}

\vspace{0.5cm}

%% --- Concepto clave ---

\begin{conceptbox}
\textbf{\textcolor{purpura}{\large Concepto clave}}\\[0.3cm]
Los posesivos en espa\~{n}ol concuerdan con el \textbf{objeto pose\'ido}, no con el poseedor.
\begin{itemize}[nosep, leftmargin=1.5em]
\item \textit{mi libro / mis libros} $\rightarrow$ cambia seg\'un el n\'umero del nombre que sigue.
\item \textit{su hermano / sus hermanos} $\rightarrow$ <<su>> no distingue si el poseedor es \'el, ella o usted.
\end{itemize}
\vspace{0.2cm}
Las formas \pos{mi}/\pos{tu}/\pos{su} son \textbf{invariables en g\'enero} (mi casa, mi libro) pero \textbf{variables en n\'umero} (mi/mis, tu/tus, su/sus).\\[0.2cm]
Las formas \pos{nuestro}/\pos{vuestro} S\'I var\'ian en \textbf{g\'enero Y n\'umero}: nuestro perro, nuestra amiga, nuestros perros, nuestras amigas.
\end{conceptbox}

\vspace{0.5cm}

%% --- Errores frecuentes por L1 ---

\subsection*{\textcolor{purpura}{Errores frecuentes por L1}}

\begin{itemize}[leftmargin=1.5em]
\item \textbf{Hablantes de ingl\'es:} <<su>> para his/her/its/your (formal)/their --- ambig\"uedad que en ingl\'es no existe. No hay error de forma, pero s\'i de comprensi\'on.
\item \textbf{Hablantes de franc\'es:} interferencia \textit{mon/ma/mes} $\rightarrow$ tienden a marcar g\'enero en mi/tu/su (\textit{*<<ma madre>>}). En espa\~{n}ol, mi/tu/su NO marcan g\'enero.
\item \textbf{Hablantes de italiano:} interferencia \textit{mio/mia} $\rightarrow$ misma tendencia.
\item \textbf{Todos:} confusi\'on singular/plural: \textit{*<<mi hermanos>>} (correcto: \textit{<<mis hermanos>>}).
\end{itemize}

\vspace{0.3cm}

%% --- Lo que el estudiante ya trae ---

\subsection*{\textcolor{purpura}{Lo que el estudiante ya trae}}

<<Su>> y <<mi>> aparecieron en los textos de Vocabulario (p.34--35): \textit{<<Su abuelo se llama Carlos>>, <<Su padre es camionero>>, <<Mi padre se llama Antonio>>, <<mis hermanos>>}. Las formas de 3.\textsuperscript{a} persona est\'an interiorizadas como exposici\'on incidental. La novedad son las formas de 1.\textsuperscript{a} y 2.\textsuperscript{a} persona plural (\pos{nuestro}, \pos{vuestro}) y la concordancia de g\'enero que solo estas presentan.

\newpage

%% ============================================================
%% PÁGINA 2: TABLA MARS EARS
%% ============================================================

\begin{center}
{\Large\textbf{\textcolor{purpura}{Correspondencia MARS EARS $\leftrightarrow$ Bloque 3}}}\\[0.3cm]
{\small La M (Modelling) ya ocurri\'o como exposici\'on incidental en Vocabulario (textos p.34--35 con <<su>>, <<mi>>, <<mis>>).\\
Ciclo abreviado: el patr\'on de mi/tu/su es transparente y el estudiante ya lo ha visto;\\
la novedad (nuestro/vuestro con concordancia de g\'enero) es deducible por contraste.}
\end{center}

\vspace{0.5cm}

{\small
\renewcommand{\arraystretch}{1.45}
\begin{longtable}{|>{\bfseries\raggedright\arraybackslash}p{1.8cm}|p{4.5cm}|p{3.8cm}|p{2.8cm}|}
\hline
\rowcolor{purpura!15}
\textbf{MARS EARS} & \textbf{Qu\'e hace el estudiante} & \textbf{Material / recurso} & \textbf{Fase} \\
\hline
\endfirsthead
\hline
\rowcolor{purpura!15}
\textbf{MARS EARS} & \textbf{Qu\'e hace el estudiante} & \textbf{Material / recurso} & \textbf{Fase} \\
\hline
\endhead

M --- Modelling & \textit{(ya ocurrido)} Ley\'o textos de p.34--35 con posesivos en contexto: \textit{<<Su abuelo se llama Carlos>>, <<Mi padre se llama Antonio>>, <<mis hermanos>>}. Comprensi\'on sin focalizaci\'on. & Textos p.34--35 & Vocabulario \\
\hline

A --- Awareness Raising & Ve frases conocidas en la p\'ildora y nota que <<su>> y <<mi>> no cambian con g\'enero. Descubre que <<nuestro/vuestro>> S\'I cambian. Categoriza en dos familias (invariables en g\'enero / con concordancia). Verifica con el cuadro del libro. & P\'ildora proyectable (diap.~1--4) + cuadro p.37 & Fase 8 \\
\hline

R --- Receptive Processing & Observa el ejemplo resuelto de la Actividad 5 y verifica que entiende el criterio de selecci\'on (persona + n\'umero). Reconoce sin producir. & Ejemplo resuelto libro p.37 act.~5 & Fase 9 (inicio) \\
\hline

S --- Structured Production & Completa huecos de la Actividad 5 con \textit{weaning off}: cuadro + tarjeta $\rightarrow$ solo cuadro $\rightarrow$ de memoria. & Libro p.37 act.~5 + tarjetas Gramatips & Fase 9 \\
\hline

E --- Expansion & Integra interrogativos (Bloque 2) + posesivos en la Actividad 6: formula preguntas a partir de respuestas, con cambio de posesivo (mi$\rightarrow$tu, nuestro$\rightarrow$vuestro). Repertorio cruzado. & Libro p.37 act.~6 & Fase 10 \\
\hline

A --- Autonomous use & Pregunta a compa\~{n}eros por nombres de familiares (Act.~7) y presenta a la clase su familia y la del compa\~{n}ero con cambio mi$\rightarrow$su (Act.~8). & Interacci\'on oral + libro p.37 act.~7--8 & Fases 11--12 \\
\hline

R --- Routinisation & Consolidaci\'on distribuida: tarea 24h, menci\'on 1 semana, integrador 4 semanas. & Cuaderno + clases futuras & Cierre de secci\'on \\
\hline

S --- Spontaneity & Uso espont\'aneo en Comunicaci\'on y Destrezas cuando describe familias sin apoyo expl\'icito. & --- & Secciones \S3.3--\S3.5 \\
\hline

\end{longtable}
}

\newpage

%% ============================================================
%% PÁGINA 3: DIAPOSITIVA 1 — LAS FLECHAS DE POSESIÓN
%% ============================================================

\begin{diapoheader}{DIAPOSITIVA 1 --- LAS FLECHAS DE POSESI\'ON}
\end{diapoheader}

\vspace{0.3cm}

\begin{marsfase}{M --- Modelling}
\end{marsfase}

\vspace{0.2cm}

\textit{Modelling puro --- el profesor lee, los estudiantes miran y escuchan.}

\vspace{0.3cm}

En pantalla aparecen 8 escenas, una a una (clic). Cada escena tiene tres elementos en l\'inea: \textbf{POSEEDOR} $\rightarrow$ posesivo (en color) $\rightarrow$ \textbf{POSE\'IDO}. La flecha conecta visualmente qui\'en posee con lo que posee/a qui\'en. El posesivo est\'a sobre la flecha en el color del poseedor. Debajo, la frase completa.

\vspace{0.5cm}

\begin{center}
\renewcommand{\arraystretch}{1.8}
\begin{tabular}{|>{\raggedleft\arraybackslash}p{3.2cm}|>{\centering\arraybackslash}p{0.5cm}|>{\centering\arraybackslash}p{2.2cm}|>{\centering\arraybackslash}p{0.5cm}|>{\raggedright\arraybackslash}p{3cm}|}
\hline
\rowcolor{purpura!10}
\textbf{POSEEDOR} & & \textbf{\textcolor{purpura}{POSESIVO}} & & \textbf{POSE\'IDO} \\
\hline

Javier (yo) & $\rightarrow$ & \pos{mi} & $\rightarrow$ & hermana \\
\multicolumn{5}{|l|}{\hspace{0.5cm}\textit{<<Mi hermana se llama Alejandra.>>}} \\
\hline

Luc\'ia (yo) & $\rightarrow$ & \pos{mis} & $\rightarrow$ & hermanos \\
\multicolumn{5}{|l|}{\hspace{0.5cm}\textit{<<Mis hermanos son tres.>>}} \\
\hline

Javier ($\rightarrow$ t\'u) & $\rightarrow$ & \pos{tu} & $\rightarrow$ & madre \\
\multicolumn{5}{|l|}{\hspace{0.5cm}\textit{<<Tu madre es enfermera, ?`verdad?>>}} \\
\hline

(hablamos de Luc\'ia) & $\rightarrow$ & \pos{su} & $\rightarrow$ & madre \\
\multicolumn{5}{|l|}{\hspace{0.5cm}\textit{<<Su madre trabaja en el hotel.>>}} \\
\hline

(los padres) & $\rightarrow$ & \pos{sus} & $\rightarrow$ & abuelos \\
\multicolumn{5}{|l|}{\hspace{0.5cm}\textit{<<Sus abuelos viven en Getafe.>>}} \\
\hline

Javier y Luc\'ia & $\rightarrow$ & \pos{nuestro} & $\rightarrow$ & instituto \\
\multicolumn{5}{|l|}{\hspace{0.5cm}\textit{<<Nuestro instituto est\'a en el barrio.>>}} \\
\hline

Javier y Luc\'ia & $\rightarrow$ & \pos{nuestra} & $\rightarrow$ & calle \\
\multicolumn{5}{|l|}{\hspace{0.5cm}\textit{<<Nuestra calle es muy tranquila.>>}} \\
\hline

profe $\rightarrow$ clase & $\rightarrow$ & \pos{vuestra} & $\rightarrow$ & profesora \\
\multicolumn{5}{|l|}{\hspace{0.5cm}\textit{<<Vuestra profesora es muy simp\'atica.>>}} \\
\hline

\end{tabular}
\end{center}

\vspace{0.4cm}

\textit{El profesor lee las 8 frases se\~{n}alando poseedor $\rightarrow$ flecha $\rightarrow$ pose\'ido. Dos pasadas. La segunda m\'as despacio, se\~{n}alando el posesivo en color. No explica nada.}

\newpage

%% ============================================================
%% PÁGINA 4: DIAPOSITIVA 2 — ¿QUÉ CAMBIA A LA VEZ?
%% ============================================================

\begin{diapoheader}{DIAPOSITIVA 2 --- ?`QU\'E CAMBIA A LA VEZ?}
\end{diapoheader}

\vspace{0.3cm}

\begin{marsfase}{A --- Awareness Raising}
\end{marsfase}

\vspace{0.2cm}

\textit{Awareness --- el estudiante nota el cambio singular/plural.}

\vspace{0.3cm}

En pantalla aparecen las mismas flechas de la Diapositiva 1 pero reorganizadas en \textbf{pares m\'inimos}: la misma persona, el mismo tipo de relaci\'on, pero una vez singular y otra plural. El contraste es visual e inmediato.

\vspace{0.5cm}

%% --- Par 1 ---
\begin{tcolorbox}[colback=white, colframe=purpura!60, title={\textbf{Par 1}}, fonttitle=\bfseries\color{purpura}, colbacktitle=purpura!10, arc=3pt, boxrule=0.8pt]
\renewcommand{\arraystretch}{1.6}
\begin{tabular}{>{\raggedleft\arraybackslash}p{3cm} c >{\centering\arraybackslash}p{1.5cm} c >{\raggedright\arraybackslash}p{3cm} p{5cm}}
Javier & $\rightarrow$ & \pos{mi} & $\rightarrow$ & padre & \textit{<<Mi padre es camionero.>>} \\
Javier & $\rightarrow$ & \pos{mis} & $\rightarrow$ & padres & \textit{<<Mis padres trabajan mucho.>>} \\
\end{tabular}
\end{tcolorbox}

\vspace{0.3cm}

%% --- Par 2 ---
\begin{tcolorbox}[colback=white, colframe=purpura!60, title={\textbf{Par 2}}, fonttitle=\bfseries\color{purpura}, colbacktitle=purpura!10, arc=3pt, boxrule=0.8pt]
\renewcommand{\arraystretch}{1.6}
\begin{tabular}{>{\raggedleft\arraybackslash}p{3cm} c >{\centering\arraybackslash}p{1.5cm} c >{\raggedright\arraybackslash}p{3cm} p{5cm}}
(Luc\'ia) & $\rightarrow$ & \pos{su} & $\rightarrow$ & madre & \textit{<<Su madre trabaja en el hotel.>>} \\
(Luc\'ia) & $\rightarrow$ & \pos{sus} & $\rightarrow$ & hermanos & \textit{<<Sus hermanos estudian en el instituto.>>} \\
\end{tabular}
\end{tcolorbox}

\vspace{0.3cm}

%% --- Par 3 ---
\begin{tcolorbox}[colback=white, colframe=purpura!60, title={\textbf{Par 3}}, fonttitle=\bfseries\color{purpura}, colbacktitle=purpura!10, arc=3pt, boxrule=0.8pt]
\renewcommand{\arraystretch}{1.6}
\begin{tabular}{>{\raggedleft\arraybackslash}p{3cm} c >{\centering\arraybackslash}p{1.8cm} c >{\raggedright\arraybackslash}p{3cm} p{4.5cm}}
Javier y Luc\'ia & $\rightarrow$ & \pos{nuestro} & $\rightarrow$ & instituto & \textit{<<Nuestro instituto est\'a cerca.>>} \\
Javier y Luc\'ia & $\rightarrow$ & \pos{nuestras} & $\rightarrow$ & mochilas & \textit{<<Nuestras mochilas est\'an aqu\'i.>>} \\
\end{tabular}
\end{tcolorbox}

\vspace{0.5cm}

\textit{El profesor lee cada par se\~{n}alando la diferencia. Tras los 3 pares, tres preguntas de awareness (aparecen una a una con clic):}

\vspace{0.3cm}

\begin{awarenessbox}
\textbf{\textcolor{purpura}{Preguntas de awareness}} (una a una, con clic)

\begin{enumerate}[leftmargin=1.5em, itemsep=0.3cm]
\item \textit{<<Mirad el par 1. ?`Qu\'e ha cambiado entre <<mi padre>> y <<mis padres>>?>>}
\item \textit{<<?>Qu\'e cambia a la vez que la -s de <<mis>>? Mirad la palabra de al lado.>>}
\item \textit{<<?`Pasa lo mismo con <<su>> y <<sus>>? Comprobadlo en el par 2.>>}
\end{enumerate}

\vspace{0.2cm}

Los estudiantes responden en parejas (20 seg). Plenaria.\\
\textbf{Respuesta esperada:} cuando lo pose\'ido es plural, el posesivo lleva \textbf{-s}. El posesivo <<copia>> el n\'umero de lo que viene despu\'es.
\end{awarenessbox}

\vspace{0.4cm}

\begin{awarenessbox}
\textbf{\textcolor{purpura}{Pregunta adicional para tensi\'on cognitiva:}}\\[0.2cm]
\textit{<<Mirad el par 3: <<nuestro>> y <<nuestras>>. ?`Solo ha cambiado el n\'umero? ?`Qu\'e m\'as ha cambiado?>>}\\[0.2cm]
\textbf{Respuesta:} Tambi\'en el g\'enero: instituto = masculino, mochilas = femenino.
\end{awarenessbox}

\newpage

%% ============================================================
%% PÁGINA 5: DIAPOSITIVA 3 — ¿TODOS FUNCIONAN IGUAL? + DESAFÍO RELÁMPAGO
%% ============================================================

\begin{diapoheader}{DIAPOSITIVA 3 --- ?`TODOS FUNCIONAN IGUAL? + \textnormal{\Large$\bigstar$} DESAF\'IO REL\'AMPAGO}
\end{diapoheader}

\vspace{0.3cm}

\begin{marsfase}{A --- Awareness + S --- Structured Production}
\end{marsfase}

\vspace{0.2cm}

\textit{Awareness con cuadro + gamificaci\'on.}

\vspace{0.5cm}

%% --- Parte 1 ---

\subsection*{\textcolor{purpura}{Parte 1 --- Comprobaci\'on con el cuadro (30 seg)}}

\begin{tcolorbox}[colback=gris, colframe=purpura!50, arc=3pt, boxrule=0.5pt]
\textit{<<Abrid el libro en la p\'agina 37. Mirad el cuadro de posesivos. Ten\'eis 30 segundos. Comprobad si lo que hab\'eis descubierto coincide con el cuadro.>>}
\end{tcolorbox}

\vspace{0.3cm}

\begin{awarenessbox}
\textbf{Pregunta en pantalla:}\\[0.2cm]
\textit{<<?`Cu\'antas formas tiene <<mi>>? ?`Y cu\'antas tiene <<nuestro>>? ?`Por qu\'e esa diferencia?>>}\\[0.2cm]
\textbf{Respuesta:} mi = 2 formas (mi/mis), nuestro = 4 formas (nuestro/nuestra/nuestros/nuestras). La diferencia: \pos{nuestro}/\pos{vuestro} cambian con g\'enero Y n\'umero; los dem\'as solo con n\'umero.
\end{awarenessbox}

\vspace{0.5cm}

%% --- Parte 2 ---

\subsection*{\textcolor{purpura}{Parte 2 --- $\bigstar$ DESAF\'IO REL\'AMPAGO: Actividad 5 (p.37)}}

\begin{gamebox}
\textbf{\textcolor{dorado!80!black}{$\bigstar$ Mec\'anica del juego}}\\[0.2cm]
Clase dividida en dos equipos. El profesor lee cada frase. Los estudiantes eligen a mano alzada (dedos: 1, 2, 3). El primero que acierta suma un punto. Ritmo r\'apido --- m\'aximo 5 segundos por frase. Cada acierto suma un punto al equipo.
\end{gamebox}

\vspace{0.4cm}

\textbf{Ejemplo resuelto} (el profesor lo hace):\\
\textit{<<Daniel estudia con \underline{\hspace{1.5cm}} primo los fines de semana.>>}\\
Opciones: su / sus / mi $\rightarrow$ \pos{su}\\
\textit{Verbaliza: Daniel = \'el $\rightarrow$ su; primo = uno $\rightarrow$ singular $\rightarrow$ su (no sus).}

\vspace{0.4cm}

\renewcommand{\arraystretch}{1.6}
\begin{tabular}{r p{8.5cm} p{4cm}}
\textbf{N.\textsuperscript{o}} & \textbf{Frase} & \textbf{Opciones} \\
\midrule
1. & \textit{?`Estudias m\'usica en \underline{\hspace{1.5cm}} instituto?} & tu / su / mi $\rightarrow$ \textbf{\pos{tu}} \\[0.1cm]
2. & \textit{Yo vivo con \underline{\hspace{1.5cm}} familia.} & mis / mi / su $\rightarrow$ \textbf{\pos{mi}} \\[0.1cm]
3. & \textit{Ellos tienen \underline{\hspace{1.5cm}} ordenador en la habitaci\'on.} & sus / su / mi $\rightarrow$ \textbf{\pos{su}} \\[0.1cm]
4. & \textit{Graciela vive con \underline{\hspace{1.5cm}} t\'ios.} & su / sus / mis $\rightarrow$ \textbf{\pos{sus}} \\[0.1cm]
5. & \textit{Yo tengo \underline{\hspace{1.5cm}} libros en la cartera.} & mi / mis / sus $\rightarrow$ \textbf{\pos{mis}} \\[0.1cm]
6. & \textit{Pedro tiene \underline{\hspace{1.5cm}} bicicleta en la calle.} & sus / su / tu $\rightarrow$ \textbf{\pos{su}} \\[0.1cm]
7. & \textit{Mi hermano y yo tenemos dos mochilas iguales. \underline{\hspace{1.5cm}} mochilas son rojas.} & Nuestro / Nuestras / Sus $\rightarrow$ \textbf{\pos{Nuestras}} \\[0.1cm]
8. & \textit{La casa de mis t\'ios es muy grande y \underline{\hspace{1.5cm}} jard\'in es muy bonito.} & sus / su / nuestro $\rightarrow$ \textbf{\pos{su}} \\[0.1cm]
9. & \textit{Mis compa\~{n}eros y yo estamos muy contentos. \underline{\hspace{1.5cm}} profesor de matem\'aticas es muy simp\'atico.} & Nuestra / Nuestro / Su $\rightarrow$ \textbf{\pos{Nuestro}} \\
\end{tabular}

\vspace{0.4cm}

\begin{gamebox}
$\bigstar$ \textit{Al completar las 9 frases, el equipo ganador recibe un aplauso. El profesor no corrige con explicaci\'on --- si fallan, se\~{n}ala qui\'en posee y qu\'e es lo pose\'ido, y otro intenta.}
\end{gamebox}

\newpage

%% ============================================================
%% PÁGINA 6: DIAPOSITIVA 4 — COMPROBAD CON EL LIBRO
%% ============================================================

\begin{diapoheader}{DIAPOSITIVA 4 --- COMPROBAD CON EL LIBRO --- Actividad 5 (p.37)}
\end{diapoheader}

\vspace{0.3cm}

\begin{marsfase}{R --- Receptive Processing + S --- Structured Production}
\end{marsfase}

\vspace{0.2cm}

\textit{Aplicaci\'on directa de lo deducido a la actividad del libro.}

\vspace{0.5cm}

%% --- Parte 1 ---

\subsection*{\textcolor{purpura}{Parte 1 --- Juntos con opciones (\'items 1--3 + ejemplo)}}

\textbf{Ejemplo resuelto} (el profesor lo lee):\\[0.2cm]
\textit{<<Daniel estudia con \pos{su} primo los fines de semana.>>}\\
\textit{Verbaliza: Daniel = \'el $\rightarrow$ su; primo = uno $\rightarrow$ singular $\rightarrow$ su (no sus).}

\vspace{0.4cm}

En pantalla aparecen los \'items 1, 2 y 3 con \textbf{dos opciones} cada uno. El profesor lee la frase y los estudiantes eligen a mano alzada:

\vspace{0.3cm}

\renewcommand{\arraystretch}{1.8}
\begin{center}
\begin{tabular}{|c p{7.5cm} p{4.5cm}|}
\hline
\rowcolor{purpura!8}
\textbf{N.\textsuperscript{o}} & \textbf{Frase} & \textbf{Opciones} \\
\hline
1. & \textit{?`Estudias m\'usica en \underline{\hspace{1.5cm}} instituto?} & \textbf{\pos{tu}} / su \\
\hline
2. & \textit{Yo vivo con \underline{\hspace{1.5cm}} familia.} & mis / \textbf{\pos{mi}} \\
\hline
3. & \textit{Ellos tienen \underline{\hspace{1.5cm}} ordenador en la habitaci\'on.} & sus / \textbf{\pos{su}} \\
\hline
\end{tabular}
\end{center}

\vspace{0.3cm}

\textit{Tras cada acierto, el profesor confirma sin explicar: solo se\~{n}ala qui\'en posee y qu\'e es lo pose\'ido.}

\vspace{0.6cm}

%% --- Parte 2 ---

\subsection*{\textcolor{purpura}{Parte 2 --- Solos: \'items 4--9}}

\begin{tcolorbox}[colback=gris, colframe=purpura!50, arc=3pt, boxrule=0.5pt]
\textit{<<Ahora completad los \'items 4 a 9 solos. !`Sin mirar el cuadro!>>}
\end{tcolorbox}

\vspace{0.3cm}

Los estudiantes completan individualmente. Estos \'items son m\'as complejos: incluyen plural (\pos{sus}, \pos{mis}), formas de \pos{nuestro/-a} con concordancia de g\'enero, y casos donde el poseedor no es evidente.

\vspace{0.3cm}

\renewcommand{\arraystretch}{1.6}
\begin{center}
\begin{tabular}{|c p{11cm}|}
\hline
\rowcolor{purpura!8}
\textbf{N.\textsuperscript{o}} & \textbf{Frase} \\
\hline
4. & \textit{Graciela vive con \underline{\hspace{2cm}} t\'ios.} \\
\hline
5. & \textit{Yo tengo \underline{\hspace{2cm}} libros en la cartera.} \\
\hline
6. & \textit{Pedro tiene \underline{\hspace{2cm}} bicicleta en la calle.} \\
\hline
7. & \textit{Mi hermano y yo tenemos dos mochilas iguales. \underline{\hspace{2cm}} mochilas son rojas.} \\
\hline
8. & \textit{La casa de mis t\'ios es muy grande y \underline{\hspace{2cm}} jard\'in es muy bonito.} \\
\hline
9. & \textit{Mis compa\~{n}eros y yo estamos muy contentos. \underline{\hspace{2cm}} profesor de matem\'aticas es muy simp\'atico.} \\
\hline
\end{tabular}
\end{center}

\vspace{0.5cm}

%% --- Pregunta interlingüística ---

\begin{awarenessbox}
\textbf{\textcolor{purpura}{Pregunta interling\"u\'istica al cerrar:}}\\[0.2cm]
\textit{<<En vuestro idioma, ?`el posesivo cambia con la persona que posee o con la cosa pose\'ida?>>}
\end{awarenessbox}

\vspace{0.4cm}

\begin{tcolorbox}[colback=dorado!10, colframe=dorado!80!black, arc=3pt, boxrule=0.8pt]
\textbf{$\bigstar$ Reparta la tarjeta Gramatips de posesivos (Caja 3).}
\end{tcolorbox}

\end{document}
